\section{A corrected uniform crown-cover lemma}
\label{sec:corrected-crowns}

This section gives an ordinary mathematical proof of the uniform
crown-cover statement used in the FLC reduction.  All Lean declarations
attached to the nodes below are planned declarations.  None of these nodes
is marked as formally proved.  The proof uses full pointed crowns, including
incomplete crowns on supertile boundaries, and replaces both problematic
estimates in the Appendix: the endpoint estimate by a minimal prototile side
length, and the sine of an exterior angle exceeding $\pi$.

Throughout, polygons are compact simple polygonal disks with nonempty
interior, positive-length edges, and carrier equal to the closure of their
interior.  They may be nonconvex.  All distances and diameters are Euclidean.
Write $B(v,r)$ for the open ball of radius $r$ about $v$.

\begin{definition}[Pointed crown catalogue and sectors]
\label{def:crown-sector}
\lean{Tilings.Crowns.pointedCatalogue}
\uses{def:polygon,def:protoset,def:placed,def:tiling,def:crown,def:flc}
A \emph{full pointed crown} at a tile vertex $v$ consists of $v$ together
with every tile whose closed carrier contains $v$.  Fix a finite catalogue
$\mathcal C$ of such pointed configurations, up to translation.  Every crown
to which the following results are applied must occur in this catalogue.
For supertiles the catalogue includes incomplete boundary crowns, rather
than only complete crowns from plane tilings.

If a polygon $A$ contains $v$, let $T_vA$ be its closed local tangent cone:
the whole plane when $v$ is an interior point, its inward halfplane when
$v$ lies strictly inside a straight boundary edge, and its interior wedge
when $v$ is a genuine corner.  Collinear listed vertices use the halfplane
case.  For a pointed crown $K$ set
\[
  \Omega(K)=\bigcup_{A\in K}T_vA,
  \qquad \operatorname{Sec}(K,r)=B(v,r)\cap\Omega(K).
\]
The union here includes tiles for which the centre $v$ lies on an edge
rather than at one of their own listed vertices.
\end{definition}

\begin{lemma}[Local polygon cones]
\label{lem:polygon-local-cone}
\lean{Tilings.Crowns.polygonLocalCone}
\uses{def:polygon,def:placed,def:crown-sector}
For every polygon $A$ and every $v\in A$, there is a number $r>0$ such that
\[
  A\cap B(v,r)=T_vA\cap B(v,r).
\]
\end{lemma}
\begin{proof}
For an interior point take a ball contained in the interior.  At a boundary
point, merge consecutive collinear boundary segments for this local
argument.  The compact union of boundary edges not incident to the point
has positive distance from it.  Decrease the radius further to avoid the
other endpoints of the one or two incident edges.  The boundary inside the
ball is consequently one straight line segment, or the two rays at a
genuine corner.  The simple polygonal disk lies locally on its prescribed
interior side, giving respectively the inward halfplane or interior wedge.
Keeping collinear vertices in the original vertex list changes neither
the disk nor this local description.
\end{proof}

\begin{lemma}[A covered tile belongs to the covering patch]
\label{lem:covered-tile-belongs}
\lean{Tilings.Patch.mem_of_carrier_subset_union}
\uses{def:polygon,def:placed}
Let $F$ be a finite patch and let $X$ be a tile whose interior is disjoint
from the interior of every tile of $F$ other than $X$ itself. If the
carrier of $X$ is contained in $\bigcup F$, then $X\in F$.
\end{lemma}
\begin{proof}
If $X\notin F$, the nonempty open set $\operatorname{int}X$ is covered by
the finitely many carriers of $F$ while meeting none of their interiors.
It is therefore covered by their finitely many boundary segments, which
have empty interior; a finite union of segments cannot contain a nonempty
open set. This is the regular-closedness step used, without being named,
in the paper's Lemma~2 (``covers and hence includes'') and in
Lemma~\ref{lem:finite-neighbourhoods}; both cite this lemma.
\end{proof}

\begin{lemma}[Uniform sector radius and vertex separation in crowns]
\label{lem:sector-radius}
\lean{Tilings.Crowns.uniformSectorRadius}
\uses{def:crown-sector,lem:polygon-local-cone}
There are $r_*>0$ and $\delta_*>0$, depending only on $\mathcal C$, such that
every represented pointed crown $K$ with centre $v$ satisfies
\[
  \bigl(\bigcup K\bigr)\cap B(v,r_*)
       =\Omega(K)\cap B(v,r_*),
\]
and every listed vertex $w\ne v$ of any tile of $K$ satisfies
$\lVert w-v\rVert\ge\delta_*$.  In particular a complete crown covers
$B(v,r_*)$.
\end{lemma}
\begin{proof}
For one crown take the minimum of the finitely many radii in
Lemma~\ref{lem:polygon-local-cone}, one for each incident tile.  This gives
the displayed equality.  Take a positive minimum of these radii over the
finite catalogue, capping them by $1$ if desired.  A largest admissible
radius need not be attained and is not used.

The collection of positive distances from the distinguished centre of a
catalogue crown to a listed vertex of one of its tiles is finite.  Its
minimum is positive; if the collection were empty, use $\delta_*=1$.
Translations preserve both conclusions.  Collinear listed vertices can be
included in this minimum.  A complete crown has cone equal to the whole
plane, so its sector is the ball.
\end{proof}

\begin{definition}[Geometric macro boundary]
\label{def:macro-boundary}
\lean{Tilings.Crowns.geometricMacroBoundary}
\uses{def:polygon,def:protoset,def:substitution}
Let $P$ be the finite prototile family and let $\theta>1$ be the expansion
factor.  A finite patch $M$ considered below has pairwise disjoint tile
interiors and union
\[
 R=u+\theta^j A,\qquad A\in P,\quad j\in\mathbb N.
\]
For geometric arguments merge consecutive collinear edges of the root
polygon $A$.  Their scaled copies are called \emph{geometric macro edges},
and their endpoints \emph{genuine macro corners}.  Every genuine corner
has interior angle different from $\pi$.  Dummy collinear vertices remain
in child-tile and crown data; only the geometric macro boundary is simplified.
\end{definition}

\begin{lemma}[Genuine macro corners are child vertices]
\label{lem:macrocorner}
\lean{Tilings.Crowns.macrocornerIsChildVertex}
\uses{def:tiling,def:macro-boundary,lem:polygon-local-cone}
Every genuine macro corner $q$ is a listed vertex of at least one child
tile of $M$.
\end{lemma}
\begin{proof}
Because $M$ is finite and its carriers are compact, a sufficiently small
ball at $q$ meets only child tiles containing $q$.  If none has a vertex
there, each incident tile has either the whole plane or a halfplane as
its local cone.  The whole-plane case would make $q$ interior to $R$.
Two open halfplanes have nonempty intersection unless they are the two
opposite halfplanes of the same line.  Such intersection would violate
disjoint child interiors.  Thus the local union is one halfplane, or the
whole plane in the opposite-halfplane case.  Neither is the nonstraight
wedge at a genuine macro corner.  This contradiction proves the assertion.

A dummy $180$-degree macro vertex need not be a child vertex.  We never
apply this assertion to such a vertex.
\end{proof}

\begin{lemma}[Corrected distance to a macro edge endpoint]
\label{lem:endpoint-gap}
\lean{Tilings.Crowns.incidentMacroEdgeVertexGap}
\uses{def:crown-sector,def:macro-boundary,lem:sector-radius,lem:macrocorner}
Assume all full crowns of $M$ are represented in $\mathcal C$.  Put
$h=\min(r_*,\delta_*)$.  If $q$ is a genuine macro corner and $v$ is a
child-tile vertex in the relative interior of an incident geometric macro
edge, then
\[
  \lVert v-q\rVert\ge h>0.
\]
\end{lemma}
\begin{proof}
Suppose $\lVert v-q\rVert<h$, and let $B$ be a child tile with vertex $v$.
By Lemma~\ref{lem:macrocorner} the crown $K_M(q)$ is represented in
$\mathcal C$.  Near $v$, both $R$ and the tangent wedge of $R$ at $q$
agree with the inward halfplane of the incident macro edge.  Choose this
neighbourhood inside $B(q,r_*)$ and away from the other ray at $q$.
Lemma~\ref{lem:sector-radius} shows that the inward half-neighbourhood is
covered by the tiles of $K_M(q)$.

Regular closedness of $B$ supplies a nonempty open subset of its interior
arbitrarily close to $v$.  It lies in that inward half-neighbourhood.  If
$B$ were not in $K_M(q)$, this open subset would be covered by the finite
union of its other tile carriers.  The finitely many boundary line
segments cannot cover a nonempty plane open set.  A remaining point would
belong to the interiors of $B$ and of a crown tile, a contradiction.
Hence $B\in K_M(q)$.  Its vertex $v\ne q$ is then among the vertices used
to define $\delta_*$, giving $\lVert v-q\rVert\ge\delta_*$, again a
contradiction.

The restriction to an incident macro edge is essential to this argument.
For a vertex on a nonincident macro edge we use the separate root
separation constant below.  Merely belonging to the closed union of a
crown does not imply that every tile containing the point is in the crown.
\end{proof}

\begin{definition}[Positive root separation constants]
\label{def:crown-separation}
\lean{Tilings.Crowns.rootSeparationConstants}
\uses{def:protoset,def:macro-boundary}
Choose $\nu>0$ bounded above by every distance from a genuine root corner
to a root edge other than its two incident edges.  Choose $\lambda>0$
bounded above by every distance between two nonadjacent geometric root
edges.  These are positive finite minima of distances between disjoint
compact sets.  If a relevant collection is empty, use $1$ for that bound.

For each reflex root corner let $\beta\in(0,\pi)$ be its exterior wedge
angle.  Define
\[
  \mu_* =\min\bigl(\{\pi/2\}\cup
          \{\beta:\text{reflex root corners}\}\bigr).
\]
Thus $0<\mu_*\le\pi/2$.  Convex exterior angles greater than $\pi$ and
dummy straight corners are excluded.  Scaling by $\theta^j\ge1$ does
not decrease the distance bounds $\nu$ and $\lambda$.
\end{definition}

\begin{lemma}[Distance to an adjacent ray]
\label{lem:ray-bound}
\lean{Tilings.Crowns.distanceToReflexBoundaryRay}
\uses{def:crown-separation}
Let two rays from $q$ enclose an angle $\beta\in(0,\pi)$, and let $v$ be
on the first ray at distance $t>0$ from $q$.  Its distance to the second
ray is $t\sin\beta$ when $\beta\le\pi/2$, and $t$ when
$\beta\ge\pi/2$.  For a reflex root angle it is therefore at least
$t\sin\mu_*$.
\end{lemma}
\begin{proof}
For unit directions $a,b$ with $a\mathbin{\cdot}b=\cos\beta$, minimize
\[
  \lVert ta-sb\rVert^2=t^2+s^2-2ts\cos\beta
\]
over $s\ge0$.  The minimum occurs at $s=t\cos\beta$ when this is
nonnegative, and at $s=0$ otherwise.  This gives the two formulas.  On
$(0,\pi/2]$ sine is increasing, and for larger angles the ray-distance
formula gives $t\ge t\sin\mu_*$.  Using the whole supporting line in
the obtuse case would give the wrong estimate.
\end{proof}

\begin{lemma}[Uniform relative crown neighbourhood]
\label{lem:relative-crown-ball}
\lean{Tilings.Crowns.uniformRelativeCrownBall}
\uses{def:crown-sector,def:macro-boundary,def:crown-separation,lem:sector-radius,lem:macrocorner,lem:endpoint-gap,lem:ray-bound}
Put
\[
 d=\tfrac12\min\{r_*,\nu,\lambda,h\sin\mu_*\}>0.
\]
For every child vertex $v$ of $M$,
\[
 R\cap B(v,d)\subseteq\Omega(K_M(v))\cap B(v,d)
                    \subseteq\bigcup K_M(v).
\]
\end{lemma}
\begin{proof}
The local tangent cone of $R$ at $v$ equals the union of the child cones
at $v$: intersect the finite patch with a sufficiently small ball and
use Lemma~\ref{lem:polygon-local-cone}.  If $v$ is interior to $R$, that
cone is the whole plane and its crown is complete.

At a genuine macro corner, a ray outside its tangent wedge is outside
$R$ initially.  A first reentry cannot meet either incident ray again,
so lies on a nonincident macro edge at distance at least
$\theta^j\nu\ge\nu>d$.  No outside-wedge point of $R$ is within $d$.

At an interior point of a geometric macro edge the cone is its inward
halfplane.  A ray into the open outward halfplane cannot reenter through
a nonadjacent macro edge before travelling $\theta^j\lambda\ge\lambda$.
An adjacent edge at a convex endpoint is wholly on the inward side, so
cannot cause reentry.  At a reflex endpoint $q$,
Lemma~\ref{lem:endpoint-gap} gives $\lVert v-q\rVert\ge h$, and
Lemma~\ref{lem:ray-bound} gives distance at least $h\sin\mu_*>d$ to
the adjacent edge.  This excludes reentry in all cases.  Finally
$d\le r_*$ and Lemma~\ref{lem:sector-radius} give the last inclusion.
\end{proof}

\begin{lemma}[A uniform ball across a trimmed tile edge]
\label{lem:trimmed-edge-ball}
\lean{Tilings.Crowns.trimmedEdgeTwoTileBall}
\uses{def:polygon,def:protoset,def:tiling,lem:polygon-local-cone,lem:relative-crown-ball}
For the number $d$ above there is $\eta>0$, depending only on $P$ and
$d$, with the following property.  If translated tiles $A,B$ have
disjoint interiors, their edges overlap near $e$, the tiles lie locally
on opposite sides there, and $e$ has distance at least $d/2$ from every
endpoint of those edges, then $B(e,\eta)\subseteq A\cup B$.
\end{lemma}
\begin{proof}
For each prototile edge take its closed trimmed segment of points at
distance at least $d/2$ from both endpoints.  Whenever nonempty, this
compact segment is disjoint from the compact union of all other boundary
edges of that polygon.  Take a positive lower bound $\eta$ for all these
finitely many distances; use $1$ if there are no nonempty trimmed
segments.  Translations preserve the bound.

Inside $B(e,\eta)$ neither tile has another boundary edge.  Its one edge
separates the ball into two connected half-balls, and the local inward
half-ball remains in its interior because no other boundary can be
crossed.  The two opposite inward half-balls, together with their common
line, cover the ball.
\end{proof}

\begin{lemma}[Exterior excursions across nearby macro edges]
\label{lem:exterior-excursion}
\lean{Tilings.Crowns.noShortExteriorExcursion}
\uses{def:macro-boundary,def:crown-separation,lem:macrocorner,lem:relative-crown-ball}
Let $H$ be convex with diameter less than
$\min\{\lambda,d\sin(\mu_*/2)\}$.  If every point of $H$ has distance
at least $d/2$ from every child vertex, then $H\cap R$ is convex.
\end{lemma}
\begin{proof}
Otherwise a segment joining two points of $H\cap R$ has an open exterior
excursion.  Let $E,F$ be its exit and reentry points on the macro boundary.
The finite polygonal boundary gives such endpoints; a segment contained
in the boundary is not an exterior excursion.  Both endpoints are in $H$.
They cannot lie on the same macro edge, since the segment between two
points of one edge lies in $R$.  Nonadjacent edges give
$\lVert E-F\rVert\ge\lambda$, contradicting the diameter bound.

For adjacent edges let $q$ be their genuine common corner.  It is a child
vertex by Lemma~\ref{lem:macrocorner}, so both distances from $E,F$ to
$q$ are at least $d/2$.  The excursion must cross a reflex exterior
wedge: at a convex corner a chord to the adjacent edge goes into the
inward halfplane at its exit point.  Writing $a=\lVert E-q\rVert$,
$b=\lVert F-q\rVert$ and using the reflex exterior angle $\beta$ gives
\[
 \lVert E-F\rVert^2=(a-b)^2+4ab\sin^2(\beta/2)
       \ge d^2\sin^2(\beta/2)
       \ge d^2\sin^2(\mu_*/2).
\]
Here $\beta/2\in(0,\pi/2)$, where sine is increasing.  Again this
contradicts the diameter bound.  Every segment in $H\cap R$ stays in $R$.
\end{proof}

\begin{lemma}[Convexification preserves the diameter bound]
\label{lem:crown-convexification}
\lean{Tilings.Crowns.closedConvexHullDiameter}
\uses{def:polygon}
If $S$ has finite diameter, its closed convex hull $H$ has diameter no
greater than that of $S$.
\end{lemma}
\begin{proof}
For finite convex combinations $x=\sum_i a_i p_i$ and
$y=\sum_j b_j q_j$ of points of $S$, use
$x-y=\sum_{i,j}a_i b_j(p_i-q_j)$ and the triangle inequality.
The coefficients sum to $1$, giving the desired distance bound.
Passing to limits gives the same bound in the closed convex hull.
\end{proof}

\begin{lemma}[Uniform crown cover for a plane tiling]
\label{lem:crown-cover-plane}
\lean{Tilings.Crowns.uniformCrownCoverPlane}
\uses{def:tiling,def:crown-sector,lem:localfinite,lem:sector-radius,lem:trimmed-edge-ball,lem:crown-convexification}
There is $\varepsilon>0$ such that, in every plane tiling by $P$ whose
full pointed crowns are represented in $\mathcal C$, every set $S$ of
diameter less than $\varepsilon$ is contained in one crown.
\end{lemma}
\begin{proof}
Use $\varepsilon<\min\{\eta,d/2\}$ and replace $S$ by its closed convex
hull $H$.  If $H$ comes within $d/2$ of a tile vertex $v$, all of $H$ is
within distance $d$ of $v$.  The complete crown at $v$ covers this ball.

Otherwise $H$ stays at distance at least $d/2$ from every vertex.  If one
tile contains $H$, take any crown containing that tile.  If not, choose
a tile $A$ meeting $H$, a point of $A\cap H$, and a point of $H\setminus A$.
Their segment exits $A$ at a nonvertex point $e$ of an edge.  Near $e$
incident tile carriers are halfplanes.  Disjoint interiors and plane
coverage force a tile $B$ on the opposite side.  The point $e$ is at
distance at least $d/2$ from both tiles' edge endpoints, so
Lemma~\ref{lem:trimmed-edge-ball} covers $H\subset B(e,\eta)$ by $A\cup B$.
Their shared edge segment has an endpoint which is a vertex of at least
one tile.  The full crown there contains both tiles and hence $H$.
An empty $S$ is covered by any crown.
\end{proof}

\begin{lemma}[Uniform relative crown cover for a supertile]
\label{lem:crown-cover-supertile}
\lean{Tilings.Crowns.uniformCrownCoverSupertile}
\uses{def:substitution,def:crown,def:flc,def:crown-sector,def:macro-boundary,lem:relative-crown-ball,lem:trimmed-edge-ball,lem:exterior-excursion,lem:crown-convexification}
There is $\varepsilon>0$ such that every finite supertile patch $M$ above,
whose full pointed crowns are represented in $\mathcal C$, has a crown
covering $S\cap R$ whenever $\operatorname{diam}(S)<\varepsilon$.
One may take
\[
 \varepsilon=\tfrac12\min
   \{\eta,d/2,\lambda,d\sin(\mu_*/2)\}>0.
\]
\end{lemma}
\begin{proof}
Replace $S$ by its closed convex hull $H$ before dividing into cases.
Lemma~\ref{lem:crown-convexification} preserves the strict diameter
bound.  If $H$ comes within $d/2$ of a child vertex $v$, every point of
$H$ is within $d$ of $v$.  Lemma~\ref{lem:relative-crown-ball} therefore
covers $H\cap R$ by $K_M(v)$.

In the remaining case $H$ stays at distance at least $d/2$ from all child
vertices.  Lemma~\ref{lem:exterior-excursion} shows that $W=H\cap R$ is
convex.  If empty, any crown suffices; if contained in one tile, any crown
containing that tile suffices.  Otherwise choose $A$ meeting $W$ and a
segment in $W$ from $A$ to a point outside $A$.  Let $e$ be an exit point.
It is not a child vertex.  A small neighbourhood meets only incident
tiles, each with a halfplane cone.  A segment leaving $A$ while staying
in $R$ requires an incident tile $B$ on the opposite side of the edge.
Indeed, with only the inward halfplane of $A$ present, such an exit is
locally impossible; a segment tangent to the edge remains in $A$ until
an endpoint, which has been excluded.

All relevant endpoints are at least $d/2$ from $e$.  Hence
Lemma~\ref{lem:trimmed-edge-ball} gives $B(e,\eta)\subset A\cup B$.
Since $e\in W$ and $\operatorname{diam}(W)<\varepsilon<\eta$, this union
covers $W$.  An endpoint of the positive-length overlap of the two edges
is a vertex of at least one of the tiles, and its full crown contains
$A\cup B$.  It therefore covers $W$ and the original $S\cap R$.
The overlap need not be an entire edge of both tiles: the argument does
not assume side-to-side matching.
\end{proof}

\begin{lemma}[Pointing a finite crown catalogue]
\label{lem:pointed-flc}
\lean{Tilings.Crowns.pointedFLCOfFLC}
\uses{def:crown,def:flc,lem:localfinite}
If the allowed full crown patches have only finitely many translation
classes, their pointed versions also have only finitely many translation
classes.
\end{lemma}
\begin{proof}
Choose one finite patch for each unpointed translation class.  A crown
centre is a listed vertex of at least one of its tiles, so each chosen
patch has only finitely many possible centres.  Retain all centre choices
which realize an allowed full crown.  This gives a finite pointed
catalogue, including incomplete crowns whenever the original catalogue
includes boundary crowns of finite supertiles.
\end{proof}

\begin{theorem}[Corrected FLC crown-cover statement]
\label{thm:corrected-crown-cover}
\lean{Tilings.Crowns.flcUniformCrownCover}
\uses{def:substitution,def:flc,lem:pointed-flc,lem:crown-cover-plane,lem:crown-cover-supertile}
For a finite polygonal substitution with the article's FLC condition,
using all allowed full crowns, there is a uniform $\varepsilon>0$ satisfying both
the plane and supertile crown-cover conclusions above.
\end{theorem}
\begin{proof}
Lemma~\ref{lem:pointed-flc} gives representatives for the finite translation
classes of allowed pointed crowns.  Every crown of the tilings in the plane conclusion is
allowed by assumption, and every full crown of a supertile is allowed
by the definition of the substitution catalogue, including those on its
boundary.  Apply the preceding two lemmas, decreasing the two positive
constants to their minimum if needed.
\end{proof}

This proves the geometric statement required by the FLC argument.  It
does not address the independent geometric grouping and routing
obligations for quasi-squares, nor the decoration repair in the
composition construction.
